\documentclass[11pt,a4paper]{article}

\usepackage[utf8]{inputenc}
\usepackage[T1]{fontenc}
\usepackage[french]{babel}
\usepackage{amsmath,amssymb,amsthm}
\usepackage{geometry}
\usepackage{hyperref}
\usepackage{url}
\usepackage{booktabs}
\usepackage{microtype}
\usepackage{pgfplots}
\pgfplotsset{compat=1.18}
\usetikzlibrary{patterns}

\geometry{margin=2.5cm}
\hypersetup{colorlinks=true, linkcolor=blue, urlcolor=blue, citecolor=blue,
            pdftitle={Une dérivation arithmétique du coefficient de Casimir},
            pdfauthor={Yan Senez}}

\newtheorem{theoreme}{Théorème}
\newtheorem{proposition}{Proposition}
\newtheorem{lemme}{Lemme}
\newtheorem{definition}{Définition}
\newtheorem{remarque}{Remarque}
\newtheorem{conjecture}{Conjecture}

\newcommand{\zetaR}{\zeta_R}
\newcommand{\actives}{\{3,5,7\}}
\newcommand{\PTset}{\{2,3,5,7\}}
\newcommand{\echoes}{\{p \geq 11\}}
\newcommand{\Nspatial}{N_{\rm spatial}}
\newcommand{\Ntrans}{N_{\rm transverse}}
\newcommand{\IPT}{I_{\rm PT}}

\title{\textbf{Une dérivation arithmétique du coefficient de Casimir}\\
       {\large depuis la Théorie de la Persistance, et une signature falsifiable des nombres premiers d'écho}}
\author{Yan Senez \\ {\small Recherche indépendante \quad \texttt{yan.senez@protonmail.com}}}
\date{Mai 2026}

\begin{document}

\maketitle

\begin{abstract}
Nous montrons que le coefficient standard de la pression de Casimir entre deux plaques parfaitement
conductrices, $P/a^{-4} = -\pi^{2}/240$, admet une dérivation arithmétique complète depuis le cadre
de la Théorie de la Persistance (PT)~\cite{Senez2026}, sans paramètre ajusté. La dérivation
factorise le coefficient en deux objets PT-natifs~: un préfacteur géométrique
$\Nspatial / (2 \cdot (2\pi)^{\Ntrans})$ identifié entièrement par le décompte dimensionnel
de la cascade des nombres premiers actifs, et un facteur eulérien
$\zeta(4) = \prod_{p}(1-p^{-4})^{-1}$ tronqué à l'ensemble premier PT-naturel
$\PTset$. La troncature reproduit le coefficient standard à $131$~ppm près~; le résidu
est rigoureusement égal à la queue eulérienne sur les nombres premiers d'écho $\echoes$. La
construction est universelle~: pour les dimensions spatiales $d \in \{1,\dots,5\}$, le résidu de
troncature égale $\prod_{p \geq 11}(1-p^{-(d+1)})^{-1}$ à la précision machine pour $d \geq 3$.
Nous formulons une conjecture quantitative falsifiable~: une mesure différentielle de Casimir à
mieux que $100$~ppm doit révéler une déviation de $1{,}31 \times 10^{-4}$ par rapport à la
prédiction QED standard. Une section finale précise ce que la lecture PT permet d'exclure
mécaniquement~: aucune extraction nette d'énergie d'une cavité Casimir statique n'est compatible
avec l'identité fondamentale du cadre.
\end{abstract}

\section{Introduction}

L'effet Casimir est la force macroscopique entre deux plaques métalliques non chargées prédite
par Casimir en 1948~\cite{Casimir1948} et confirmée expérimentalement à mieux que le pour-cent par
Lamoreaux~\cite{Lamoreaux1997}, Mohideen et Roy~\cite{Mohideen1998},
Bressi \emph{et al.}~\cite{Bressi2002}, et Decca \emph{et al.}~\cite{Decca2003,Decca2007}. Dans la
limite plaques idéales (parfaitement conductrices, température nulle, sans dispersion matérielle),
la pression prédite est
\begin{equation}
P_{\rm Casimir}(a) = -\frac{\pi^{2}}{240\,a^{4}},
\label{eq:standard-coefficient}
\end{equation}
où $a$ est la séparation entre plaques en unités naturelles $\hbar = c = 1$. Le coefficient
$\pi^{2}/240$ provient traditionnellement de la sommation régularisée des modes de point zéro
\emph{via} $\zetaR(-3) = 1/120$~\cite{Mostepanenko2009,Milton2001}.

Cette dérivation est interprétée habituellement comme la manifestation du vide quantique~: une
somme sur les modes dont la différence entre configurations bornée et libre donne une force finie
attractive. La somme infinie d'énergie de point zéro est traitée par continuation analytique ou
régularisation zêta, méthodes qui présupposent un spectre continu.

La Théorie de la Persistance (PT)~\cite{Senez2026} reformule ce cadre sans recours au vide
quantique. En PT, chaque observable du Modèle Standard est dérivé d'un unique axiome arithmétique
$s = 1/2$ et du point fixe structurel $\mu^{*} = 15$ du crible d'Ératosthène. Les
$43$~observables du Modèle Standard sont reproduits depuis les seuls
$\{\gamma_{p}, \sin^{2}\theta_{p}, \mu^{*}\}$ avec une précision médiane de $0{,}06\%$ et zéro
paramètre ajusté.

Une question naturelle se pose~: le coefficient \eqref{eq:standard-coefficient} admet-il la même
dérivation~? Nous répondons positivement. Le coefficient se factorise en un préfacteur géométrique
identifié en termes PT et un facteur eulérien tronqué à l'ensemble premier PT-naturel. Le résidu
de troncature est exactement la queue eulérienne sur les primes d'écho, fournissant une
signature falsifiable.

L'article est structuré comme suit. La Section~\ref{sec:standard} rappelle la dérivation standard
en dimension quelconque. La Section~\ref{sec:pt-decomposition} introduit la classification PT des
nombres premiers. La Section~\ref{sec:c3} démontre le résultat principal de troncature
(Proposition~\ref{prop:c3-residue}) et l'unicité PT-naturelle (Théorème~\ref{thm:c3-uniqueness}). La Section~\ref{sec:prefactor} identifie le préfacteur
géométrique en termes PT-natifs. La Section~\ref{sec:universality} étend le résultat à des
dimensions arbitraires. La Section~\ref{sec:falsifiability} formule la conjecture falsifiable et
propose un protocole expérimental. La Section~\ref{sec:no-extraction} précise pourquoi la lecture
PT exclut toute extraction d'énergie nette d'une cavité Casimir statique. Une discussion
(Section~\ref{sec:discussion}) replace ces résultats dans le programme PT.

\section{La dérivation standard~: le squelette analytique}
\label{sec:standard}

La pression de Casimir idéale entre deux plaques parfaitement conductrices séparées de $a$ en
$d+1$~dimensions d'espace-temps ($d$~spatiales), pour un champ massless avec $\nu$ modes
indépendants par fréquence et conditions de bord de Dirichlet, est~\cite{Mostepanenko2009,Klimchitskaya2009}
\begin{equation}
|P_{d}^{(\nu)}| = \nu \cdot \frac{d \cdot \Gamma\bigl((d+1)/2\bigr) \cdot \zeta(d+1)}{2 \cdot 2^{d}\,\pi^{(d+1)/2}\,a^{d+1}}.
\label{eq:dim-d-coefficient}
\end{equation}
Le facteur $\nu$ compte les modes indépendants~: $\nu = 1$ pour un scalaire simple, $\nu = 2$ pour
le champ électromagnétique en $d = 3$ (polarisations TE et TM). Le facteur $1/2$ est la
demi-différence de Casimir entre spectres bornés et libres (ce \emph{n'est pas} une normalisation
de point zéro dans la lecture PT~; voir Section~\ref{sec:no-extraction}). Le compte de modes $\nu$
en dimension supérieure dépend de la convention et n'affecte pas l'analyse de troncature C3
ci-dessous.

Pour $d = 3$, $\nu = 2$ redonne~\eqref{eq:standard-coefficient}. Pour $d = 1$, $\nu = 1$ donne
$|P_1| = \pi/(24\,a^2)$. Pour présenter un tableau dimensionnel uniforme en
Section~\ref{sec:universality}, nous adoptons la convention $\nu = 2$ partout (multiplicité
généralisée ``EM-like''), ce qui ne modifie aucun énoncé PT~: le facteur de multiplicité s'annule
dans le résidu relatif~\eqref{eq:c3-residue}.

Le coefficient \eqref{eq:dim-d-coefficient} contient deux structures distinctes~:
\begin{enumerate}
  \item Un \emph{préfacteur géométrique}
  $\mathcal{G}_{d} = d \cdot \Gamma((d+1)/2) / (2^{d}\,\pi^{(d+1)/2})$.
  \item Un \emph{facteur arithmétique} $\zeta(d+1)$.
\end{enumerate}

Le facteur arithmétique admet le développement eulérien
\begin{equation}
\zeta(s) = \prod_{p \text{ premier}} \frac{1}{1 - p^{-s}}.
\label{eq:euler-product}
\end{equation}
C'est le point d'entrée de l'analyse PT~: le coefficient \eqref{eq:dim-d-coefficient} est
structurellement un produit sur tous les premiers, modulé par un préfacteur géométrique fini.

\section{Décomposition PT du crible}
\label{sec:pt-decomposition}

PT classifie les nombres premiers selon leur rôle dans la cascade du crible
$\mu_{k+1} = \sum \{p : \gamma_{p}(\mu_{k}) > 1/2 \}$. L'unique point fixe auto-cohérent est
$\mu^{*} = 3 + 5 + 7 = 15$, avec la structure suivante~\cite{Senez2026}~:

\begin{itemize}
  \item $p = 2$~: \emph{canal de parité}. Cinématique uniquement ($T_{2} = (1)$). Porte la symétrie
  binaire $s = 1/2$.
  \item $p \in \actives$~: \emph{nombres premiers actifs dimensionnels}. $\gamma_{p}(\mu^{*}) > 1/2$
  pour $p \in \actives$, contribuant aux trois dimensions spatiales de $\mathrm{SU}(N_{c}=3)$ et
  aux observables actives.
  \item $p \in \echoes$~: \emph{nombres premiers d'écho}. $\gamma_{p}(\mu^{*}) < 1/2$.
  Contribuent aux corrections virtuelles (\emph{e.g.} VP d'écho dans $\alpha_{\rm EM}$) et aux
  observables du secteur sombre.
\end{itemize}

L'ensemble premier PT-naturel est donc $\PTset$. Nous décomposons le produit eulérien
\eqref{eq:euler-product} le long de cette classification~:
\begin{equation}
\zeta(s) = \zeta_{\rm parit\acute{e}}(s) \cdot \zeta_{\rm actifs}(s) \cdot \zeta_{\rm echo}(s),
\label{eq:pt-split}
\end{equation}
où $\zeta_{\rm parit\acute{e}}(s) = (1 - 2^{-s})^{-1}$,
$\zeta_{\rm actifs}(s) = \prod_{p \in \actives}(1-p^{-s})^{-1}$, et
$\zeta_{\rm echo}(s) = \prod_{p \geq 11}(1-p^{-s})^{-1}$.

\section{La troncature C3~: prédiction sans paramètre}
\label{sec:c3}

\begin{definition}[Troncature C3 PT-naturelle]
On définit la fonction zêta PT-tronquée
\begin{equation}
\zeta_{\rm PT}(s) := \zeta_{\rm parit\acute{e}}(s) \cdot \zeta_{\rm actifs}(s)
                   = \prod_{p \in \PTset}\frac{1}{1 - p^{-s}}.
\end{equation}
Le coefficient de pression de Casimir PT-natif C3 est
\begin{equation}
|P_{d}^{\rm PT}| := \mathcal{G}_{d} \cdot \zeta_{\rm PT}(d+1) \cdot a^{-(d+1)}.
\end{equation}
\end{definition}

\begin{proposition}[Résidu de troncature C3]
\label{prop:c3-residue}
L'erreur relative de la troncature C3 est exactement la correction multiplicative de queue
d'écho~:
\begin{equation}
\frac{|P_{d}^{(\nu)}| - |P_{d}^{{\rm PT},(\nu)}|}{|P_{d}^{(\nu)}|}
\;=\; 1 - \frac{1}{\zeta_{\rm echo}(d+1)}
\;=\; 1 - \prod_{p \geq 11}\bigl(1 - p^{-(d+1)}\bigr).
\label{eq:c3-residue}
\end{equation}
\end{proposition}

\begin{proof}
Par \eqref{eq:pt-split}, $\zeta(d+1) = \zeta_{\rm PT}(d+1) \cdot \zeta_{\rm echo}(d+1)$. En
substituant dans le coefficient standard~\eqref{eq:dim-d-coefficient}~:
\[
|P_{d}^{(\nu)}| = \nu\,\mathcal{G}_{d} \cdot \zeta_{\rm PT}(d+1) \cdot \zeta_{\rm echo}(d+1) \cdot a^{-(d+1)} = |P_{d}^{{\rm PT},(\nu)}| \cdot \zeta_{\rm echo}(d+1).
\]
Donc $|P_{d}^{{\rm PT},(\nu)}|/|P_{d}^{(\nu)}| = 1/\zeta_{\rm echo}(d+1) = \prod_{p \geq 11}(1-p^{-(d+1)})$, et
l'erreur relative \eqref{eq:c3-residue} en découle. La multiplicité $\nu$ se simplifie.
\end{proof}

La Proposition~\ref{prop:c3-residue} est une identité algébrique. Le contenu substantif de cet
article est plutôt l'énoncé structurel suivant, qui justifie pourquoi $\PTset$ est l'\emph{unique}
troncature PT-naturelle~:

\begin{theoreme}[Unicité de l'ensemble premier C3]
\label{thm:c3-uniqueness}
Soit $S \subset \mathbb{N}$ un ensemble fini de premiers tel que le coefficient tronqué
$\mathcal{G}_{d} \cdot \prod_{p \in S}(1-p^{-(d+1)})^{-1}$ approxime le coefficient
standard~\eqref{eq:dim-d-coefficient} avec une erreur relative décroissant en
$\sum_{p \notin S} p^{-(d+1)}$ uniformément pour $d \in \{1,\ldots,5\}$. Alors $S = \PTset$ est
l'unique tel ensemble engendré par la classification PT des nombres premiers
de~\cite{Senez2026} (Théorèmes T5 et T6).
\end{theoreme}

\begin{proof}[Esquisse]
La classification PT (\cite{Senez2026}, T5) sépare les premiers en $\{p : \gamma_p(\mu^*) > 1/2\}
= \actives$ (actifs dimensionnels), $\{2\}$ (parité, cinématique), et $\{p \geq 11\}$ (primes d'écho).
T6 prouve que $\gamma_p$ est strictement monotone décroissant en $p$ à $\mu^* = 15$ (preuve
algébrique avec rationnel exact $q_+ = 13/15$), donc l'ensemble actif est exactement
$\{3, 5, 7\}$. Le premier de parité $p = 2$ entre dans la somme spectrale cinématiquement (il
détermine l'existence des modes, non leur contribution dimensionnelle). Donc la troncature
PT-naturelle ``physiquement présente'' est $\PTset$, et tout ensemble plus grand incorpore des
primes d'écho que PT classifie comme non-dimensionnels.

Réciproquement, tout sous-ensemble strict $S \subsetneq \PTset$ soit omet la parité (et le
résidu acquiert un facteur $1/15$ supplémentaire), soit omet un premier actif (le résidu
acquiert $\sim 1/p^{d+1}$ à $d$ fixé, brisant le scaling uniforme en~$d$).
\end{proof}

\begin{remarque}
Le Théorème~\ref{thm:c3-uniqueness} est le résultat structurel substantif. L'accord numérique
de la Section~\ref{sec:c3} en découle alors par application de la
Proposition~\ref{prop:c3-residue} à l'unique troncature PT-naturelle $S = \PTset$.
\end{remarque}

\subsection{Évaluation numérique en 3D}

Pour $d = 3$, on calcule
\begin{align*}
\zeta_{\rm parit\acute{e}}(4) &= (1 - 2^{-4})^{-1} = 16/15,\\
\zeta_{\rm actifs}(4) &= \tfrac{81}{80} \cdot \tfrac{625}{624} \cdot \tfrac{2401}{2400} \approx 1{,}014542,\\
\zeta_{\rm PT}(4) &= 16/15 \cdot 81/80 \cdot 625/624 \cdot 2401/2400 \approx 1{,}082091,\\
\zeta(4) &= \pi^{4}/90 \approx 1{,}082323.
\end{align*}
Le préfacteur géométrique en $d = 3$ vaut
$\mathcal{G}_{3} = 3 \Gamma(2)/(8\pi^{2}) = 3/(8\pi^{2})$. Donc
\[
|P_{3}^{\rm PT}| = \frac{3}{8\pi^{2}} \cdot 1{,}082091 \cdot a^{-4}
\;\approx\; 0{,}041118 \cdot a^{-4},
\]
contre la valeur standard $|P_{3}| = \pi^{2}/240 \cdot a^{-4} \approx 0{,}041123 \cdot a^{-4}$.

L'erreur relative est de $131$~ppm, identique à $1/\zeta_{\rm echo}(4) - 1 = -\sum_{p \geq 11}\log(1-p^{-4}) \approx 1{,}31 \times 10^{-4}$, exactement comme prédit par la Proposition~\ref{prop:c3-residue}.

\section{Le préfacteur géométrique en termes PT}
\label{sec:prefactor}

Le préfacteur $\mathcal{G}_{3} = 3/(8\pi^{2})$ admet une identification complète en termes
PT-natifs~:
\begin{equation}
\frac{3}{8\pi^{2}} \;=\; \frac{\Nspatial}{2 \cdot (2\pi)^{\Ntrans}},
\label{eq:prefactor-pt}
\end{equation}
avec $\Nspatial = 3$ (le nombre de premiers actifs, équivalent au nombre de dimensions spatiales~;
voir~\cite{Senez2026}, Théorème T5) et $\Ntrans = 2$ (le nombre de directions de modes transverses
aux plaques). Le facteur $1/2$ est la normalisation standard de l'énergie de point zéro, et
$(2\pi)^{\Ntrans}$ correspond à deux copies du périmètre cyclique de $S^{1}$, limite continue de
$\mathbb{Z}/p\mathbb{Z}$ dans la construction PT.

L'identification \eqref{eq:prefactor-pt} est exacte et ne contient aucun apport analytique
au-delà de la structure PT. Le coefficient de Casimir en 3D est donc entièrement PT-natif modulo
le produit eulérien universel $\zeta(4)$, lui-même décomposé selon la classification PT.

\section{Universalité dimensionnelle}
\label{sec:universality}

Nous testons la Proposition~\ref{prop:c3-residue} pour les dimensions $d = 1$ à $d = 5$ \emph{via}
la formule~\eqref{eq:dim-d-coefficient}.

\begin{table}[ht]
\centering
\caption{Troncature C3 selon la dimension. Le résidu de troncature égale la correction
multiplicative de queue d'écho à la précision machine pour $d \geq 3$ et avec une précision
limitée par l'expansion logarithmique pour $d = 1, 2$.}
\label{tab:dimensional-scan}
\begin{tabular}{cccccc}
\toprule
$d$ & $|P_{d}| \cdot a^{d+1}$ standard & $|P_{d}^{\rm PT}| \cdot a^{d+1}$ C3 & Résidu (ppm) & Queue d'écho (ppm) & Accord \\
\midrule
$1$ & $0{,}26180$ & $0{,}25386$ & $30\,325$ & $31\,272$ & $\sim$3\% \\
$2$ & $0{,}09566$ & $0{,}09548$ & $1\,809$ & $1\,812$ & $0{,}2\%$ \\
$3$ & $0{,}04112$ & $0{,}04112$ & $131{,}0$ & $131{,}0$ & exact \\
$4$ & $0{,}01970$ & $0{,}01970$ & $10{,}3$ & $10{,}3$ & exact \\
$5$ & $0{,}01025$ & $0{,}01025$ & $0{,}84$ & $0{,}84$ & exact \\
\bottomrule
\end{tabular}
\end{table}

La pureté de la prédiction croît exponentiellement avec~$d$. C'est une conséquence quantitative
de la Proposition~\ref{prop:c3-residue}~: la correction de queue d'écho est bornée par
$\sum_{p \geq 11} p^{-(d+1)}$, qui converge super-exponentiellement à mesure que $d$ croît.

\begin{figure}[ht]
\centering
\begin{tikzpicture}
\begin{semilogyaxis}[
    width=12cm, height=7cm,
    xlabel={dimension spatiale $d$},
    ylabel={résidu (ppm)},
    xmin=0.5, xmax=5.5,
    ymin=0.3, ymax=100000,
    xtick={1,2,3,4,5},
    grid=major,
    legend pos=north east,
    legend cell align={left},
]
\addplot[thick, blue, mark=*, mark size=3pt] coordinates {
    (1, 30325) (2, 1809) (3, 131.0) (4, 10.3) (5, 0.84)
};
\addlegendentry{Résidu de troncature C3}
\addplot[thick, red, dashed, mark=square*, mark size=3pt] coordinates {
    (1, 31272) (2, 1812) (3, 131.0) (4, 10.3) (5, 0.84)
};
\addlegendentry{Queue d'écho $\prod_{p\geq11}(1-p^{-(d+1)})^{-1}-1$}
\addplot[thick, green!60!black, dotted] coordinates {
    (1, 90909) (2, 8264) (3, 751) (4, 68.3) (5, 6.21)
};
\addlegendentry{Borne $\sum_{p \geq 11} p^{-(d+1)}$}
\node[anchor=south west, fill=yellow!30, inner sep=2pt] at (axis cs:2.7,200) {point d'équilibre expérimental};
\draw[thick, ->, gray] (axis cs:3.05,300) -- (axis cs:3,140);
\end{semilogyaxis}
\end{tikzpicture}
\caption{Résidu de troncature C3 (en ppm) selon la dimension spatiale $d$. Trait bleu plein~:
résidu de troncature observé $1 - |P_d^{\rm PT}|/|P_d|$. Trait rouge tireté~: correction
multiplicative de queue d'écho prédite $\zeta_{\rm echo}(d+1) - 1$. Les deux courbes coïncident
à la précision machine pour $d \geq 3$. Trait vert pointillé~: la borne supérieure simple de
la Proposition~\ref{prop:asymptotic}. Le point 3D ($131$~ppm) est au point d'équilibre
expérimental~: suffisamment petit pour constituer une prédiction propre, suffisamment grand
pour être testable par les techniques différentielles actuelles.}
\label{fig:c3-vs-d}
\end{figure}

\begin{proposition}[Pureté asymptotique]
\label{prop:asymptotic}
La prédiction C3 approche le coefficient exact à la vitesse
\begin{equation}
\frac{|P_{d}^{(\nu)}| - |P_{d}^{{\rm PT},(\nu)}|}{|P_{d}^{(\nu)}|} \;\leq\; \sum_{p \geq 11} p^{-(d+1)} \;\sim\; 11^{-(d+1)}\;\text{quand}\;d \to \infty.
\end{equation}
\end{proposition}

La dimension physique $d = 3$ se trouve dans le \emph{point d'équilibre expérimental}~: le résidu
de $131$~ppm est suffisamment petit pour constituer une prédiction propre, mais suffisamment grand
pour être testable avec les techniques de mesure différentielle actuelles (cf.
Section~\ref{sec:falsifiability}).

\section{Falsifiabilité et signature expérimentale}
\label{sec:falsifiability}

Le résidu 3D de $131$~ppm est dans le domaine de précision atteignable par les meilleures mesures
de Casimir contemporaines. Decca \emph{et al.}~\cite{Decca2007,Decca2003} atteignent
$\lesssim 1\%$ de précision sur géométries plaque-sphère, tandis que les mesures différentielles
avec configurations conçues pour supprimer les corrections de matériau pourraient en principe
atteindre le régime sub-millième.

\begin{conjecture}[Signature PT dans Casimir]
\label{conj:pt-signature}
Une mesure de Casimir à séparation $a$ sur géométrie idéale, après soustraction de toutes les
corrections de Lifshitz et de conductivité finie connues~\cite{Lifshitz1956}, dévie de la
prédiction standard \eqref{eq:standard-coefficient} d'exactement
\begin{equation}
\Delta P / P \;=\; 1 - \prod_{p \geq 11}(1 - p^{-4}) \;=\; 1{,}31 \times 10^{-4}.
\label{eq:pt-signature}
\end{equation}
\end{conjecture}

Cette déviation \emph{n'est pas} prédite par la QED standard, qui utilise la continuation
analytique de $\zetaR$ et somme implicitement tous les premiers. La
Conjecture~\ref{conj:pt-signature} est donc une prédiction \emph{discriminante}~: la confirmation
de la déviation soutient la classification PT des primes d'écho~; un résultat nul à mieux que
$100$~ppm la falsifie.

\subsection{Protocole expérimental}

La signature \eqref{eq:pt-signature} se situe sous les incertitudes actuelles des corrections
de Lifshitz sur les mesures de pression absolue ($\gtrsim 1\%$ sur la plupart des géométries
en raison des ambiguïtés Drude vs.\ plasma, conductivité finie et corrections
thermiques~\cite{Klimchitskaya2009}). Un test direct sur le coefficient absolu n'est donc pas
actuellement faisable. La conjecture reste néanmoins falsifiable \emph{via} des schémas
\emph{différentiels} qui annulent les corrections matériau communes et isolent le facteur
arithmétique $\zeta(d+1)$. Une falsification pratique nécessite~:

\begin{enumerate}
  \item Une précision sub-$100$~ppm sur le rapport \emph{différentiel} $|P|_A/|P|_B$ entre
    deux configurations $A$ et $B$ partageant les propriétés matérielles mais différant
    uniquement dans $\mathcal{G}_d$ ou le comportement de troncature $\zeta$.
  \item Une calibration indépendante de la géométrie relative à mieux que $50$~ppm (la
    composante PT-native est fixée par le décompte dimensionnel, donc inchangée entre $A$ et $B$).
  \item Annulation des corrections de conductivité finie, rugosité et température \emph{via}
    le rapport différentiel (ces corrections sont en mode commun et disparaissent à l'ordre
    dominant).
\end{enumerate}

Un schéma différentiel concret compare la pression entre plaque-sphère (géométrie Decca
standard) et plaque-cylindre (ou plaque-cône) à séparation $a$ et matériau identiques. Les
préfacteurs géométriques diffèrent par une fonction de courbure connue~\cite{Mostepanenko2009},
tandis que le facteur arithmétique $\zeta(d+1)$ est identique~: toute déviation du rapport
mesuré par rapport à la prédiction QED au-delà de $100$~ppm teste le résidu de troncature PT.
Un schéma alternatif utilise deux matériaux aux paramètres de Drude très proches (par exemple
revêtements Au et Cu) tels que les corrections Lifshitz s'annulent à mieux que $50$~ppm dans
le rapport, tandis que la signature d'écho arithmétique est universelle.

L'état de l'art actuel (Decca \emph{et al.}~\cite{Decca2007}) atteint $\sim 0{,}2\%$ sur la
pression absolue mais est dominé par les ambiguïtés des modèles de Lifshitz. Un dispositif
différentiel conçu pour supprimer ces ambiguïtés est la cible naturelle.

\section{Pourquoi la lecture PT exclut l'extraction d'énergie}
\label{sec:no-extraction}

La popularité médiatique récente de propositions visant à extraire de l'énergie nette de cavités
Casimir statiques (dispositifs de type \emph{quantum ratchet} en cavité, brevets sur
\emph{vacuum-to-current converters}, etc.) justifie une clarification spécifique. La lecture PT
fournit un cadre direct pour qualifier ces propositions sans calcul détaillé.

\subsection{L'identité fondamentale}

Dans la formulation PT, la quantité physique observable est la divergence de Kullback-Leibler
\begin{equation}
\IPT(a, \varepsilon, \mathcal{G}) := D_{\rm KL}\bigl(P_{\rm bord}(a,\varepsilon,\mathcal{G})\,\|\,P_{\rm libre}\bigr),
\label{eq:I-PT}
\end{equation}
où $a$ est la séparation, $\varepsilon$ la fonction diélectrique, $\mathcal{G}$ la géométrie, et
$P_{\rm bord}$, $P_{\rm libre}$ sont les distributions sur les survivants du crible CRT
contraint par les bords. La pression observée est~$P = -\partial_{a}(\IPT/a^{d-1})$.

Le résultat clé est que $\IPT$ est une \emph{fonction d'état} au sens thermodynamique~: sa
valeur ne dépend que de la configuration instantanée, jamais de la trajectoire suivie pour
y arriver. C'est garanti par l'identité fondamentale du gap (GFT, monographie ch.~4
de~\cite{Senez2026})~:
\begin{equation}
\log_{2}(m) = D_{\rm KL}(P\,\|\,U_{m}) + H(P),
\label{eq:gft}
\end{equation}
qui est une identité algébrique exacte (et non une approximation thermodynamique).

\subsection{Conséquence~: pas de cycle ouvert}

Pour tout cycle fermé dans l'espace des configurations $(a, \varepsilon, \mathcal{G})$, on a
\begin{equation}
\oint \mathrm{d}\IPT = 0.
\label{eq:closed-cycle}
\end{equation}
Donc le travail mécanique net extrait sur un cycle complet est exactement nul. Ce n'est pas une
inégalité (premier principe ordinaire) mais une \emph{égalité algébrique} (anti-double-comptage).

\subsection{Conséquence pratique}

Cette propriété disqualifie immédiatement plusieurs classes de propositions d'extraction~:

\begin{enumerate}
  \item \emph{Cavités Casimir statiques avec ratchet électronique}~: la cavité étant fixe en
  $(a, \varepsilon, \mathcal{G})$, le système atteint l'équilibre détaillé. Aucun courant net
  DC ne peut persister sans injection externe. Le ratchet purement géométrique reproduit
  l'erreur de Smoluchowski-Feynman~\cite{Feynman1963}~: l'asymétrie thermodynamique apparente
  entre paroi et antenne est compensée par le bilan détaillé des taux de tunneling.

  \item \emph{Cycles à modulation de matériau} (\emph{e.g.}, schémas de type
  Pinto~\cite{Pinto1999})~: le travail nécessaire pour modifier $\varepsilon(\omega)$
  excède toujours le gain Casimir, car la modulation est elle-même un chemin dans
  $(\varepsilon)$ qui contribue à l'intégrale~\eqref{eq:closed-cycle}.

  \item \emph{Conversion vide-courant en régime stationnaire}~: l'état stationnaire impose
  $\partial_{t} \IPT = 0$, donc aucun flux d'information persistante ne traverse le système.
  La construction PT n'autorise pas de stationnarité avec extraction.
\end{enumerate}

Le seul mécanisme de \emph{conversion} (et non de création) compatible avec PT est l'effet
Casimir dynamique (DCE) observé par Wilson \emph{et al.}~\cite{Wilson2011}~: la modulation
mécanique rapide d'une plaque convertit des modes virtuels en photons réels. Le bilan
énergétique est strict~: $\hbar\omega_{\rm photon} \leq E_{\rm m\acute{e}canique\,inject\acute{e}e}$
à chaque cycle.

\subsection{Le faux espoir architectural}

Un dernier scénario théorique mérite mention. Si l'on pouvait modifier localement $\mu^{*}$
(le point fixe du crible, normalement $\mu^{*} = 15$), on changerait la pression de Casimir et
pourrait \emph{a priori} produire $\oint \mathrm{d}\IPT \neq 0$. Mais $\mu^{*} = 15$ est fixé
par auto-cohérence arithmétique (Théorème T5 de~\cite{Senez2026})~: ce n'est pas un paramètre
dynamique mais l'unique point fixe stable du flot du crible. Le modifier reviendrait à modifier
l'arithmétique elle-même, ce qui est par construction impossible dans PT.

\subsection{Synthèse}

La lecture PT ne libère pas l'effet Casimir des contraintes thermodynamiques~: elle les rend
plus rigoureuses en montrant qu'elles sont de nature arithmétique, non seulement physique.
Casimir reste une \emph{sonde}, jamais une \emph{source}.

Mais c'est une sonde exceptionnellement sensible au sens informationnel~: les $131$~ppm de
signature d'écho constituent la trace observable la plus directe que la classification PT
des nombres premiers possède une réalité structurelle au-delà du Modèle Standard. Le gain
expérimental est d'ordre informationnel, non énergétique.

\section{Discussion}
\label{sec:discussion}

La Conjecture~\ref{conj:pt-signature} reformule l'effet Casimir dans l'ontologie PT. Trois
observations structurelles en découlent.

\subsection{Casimir comme sonde du crible, non du vide}

L'effet Casimir n'est pas une manifestation du vide quantique mais une mesure du déficit
spectral du crible CRT contraint. Aucune somme infinie d'énergie de point zéro n'est nécessaire,
aucune régularisation requise~: le produit eulérien dans la décomposition PT est fini à chaque
étape. La continuation analytique standard $\zetaR(s)$ devient un outil de calcul commode
plutôt qu'un mystère ontologique.

\subsection{Les nombres premiers d'écho deviennent observables}

La queue d'écho $\prod_{p \geq 11}(1-p^{-(d+1)})^{-1}$ est, par la Proposition~\ref{prop:c3-residue},
le contenu intégral de la déviation entre la troncature PT-native et le coefficient standard.
Toute mesure de Casimir sub-pour-mille teste donc directement la classification PT des nombres
premiers en actifs, parité et d'écho. Cela connecte l'effet Casimir à l'identification du
secteur sombre cosmologique développée dans~\cite{Senez2026}, où les primes d'écho
représentent la fraction de matière noire $F_{\rm echo} \approx 95\%$ à grand $N$.

L'identification est cohérente~: ce ne sont pas \emph{deux} structures distinctes (matière noire
\emph{et} écho Casimir) mais \emph{une seule}, observée à deux échelles différentes
(cosmologique \emph{vs} laboratoire).

\subsection{Le préfacteur se réduit au comptage dimensionnel}

Le facteur $\mathcal{G}_{3} = 3/(8\pi^{2})$ ne contient aucun paramètre libre une fois la
classification PT établie ($\Nspatial = 3$, $\Ntrans = 2$). Le $\pi^{2}$ provient de la limite
$\mathbb{Z}/p\mathbb{Z} \to S^{1}$ du cycle discret, intrinsèque à la construction PT. Aucun
ajustement n'est laissé.

\subsection{Comparaison avec $\alpha_{\rm EM}$}

Une décomposition PT similaire a été démontrée pour la constante de structure
fine~\cite{Senez2026}~: $\alpha_{\rm EM}^{-1} = 137{,}0359990834$ depuis un produit en cascade
PT sur les actifs $\actives$, avec corrections VP d'échos provenant de $\{11, 13\}$. La
structure
\begin{equation}
\alpha_{\rm EM}^{-1}_{\rm habill\acute{e}} = \alpha_{\nu}^{-1} + \text{corrections d'écho}
\end{equation}
reflète la structure Casimir
\begin{equation}
P_{\rm exact} = P_{\rm PT} \cdot \zeta_{\rm echo}(d+1).
\end{equation}
Les deux observables sont construites depuis un cœur PT-actif augmenté de corrections
d'écho. La magnitude relative de la correction d'écho est beaucoup plus faible pour Casimir
($\sim 10^{-4}$ en 3D) que pour $\alpha_{\rm EM}$ ($\sim 10^{-3}$) parce que Casimir vit dans le
$\zeta(4)$ rapidement convergent plutôt que dans le produit en cascade lentement convergent de
$\alpha_{\rm EM}$.

\subsection{Position dans le programme PT}

Cet article complète une trajectoire commencée par la dérivation arithmétique d'$\alpha_{\rm EM}$
et étendue par les chantiers \texttt{PT\_RH}~\cite{Senez2026RH} sur la fonction zêta de Riemann
et \texttt{PT\_HL}~\cite{Senez2026HL} sur les conjectures de Hardy-Littlewood. Casimir occupe
une position particulière. D'autres prédictions PT testables à l'échelle du laboratoire existent
(notamment $a_\mu$ et $a_e$ \emph{via} les corrections VP d'échos au running
d'$\alpha_{\rm EM}$~\cite{Senez2026}, toutes deux compatibles avec les mesures actuelles).
Ce qui distingue Casimir est qu'il s'agit du test PT le plus \emph{direct et accessible}~: la
signature d'écho est isolée arithmétiquement (non noyée dans le running QED), le préfacteur
se réduit à du décompte dimensionnel (aucun apport analytique), et la prédiction se situe à
un niveau de précision atteignable par des mesures différentielles dédiées dans les
laboratoires universitaires actuels. Les autres prédictions prospectives (cosmologie~:
testable par Euclid 2028, DESI 2030, JUNO 2027~; négatives~: absence d'axion, de SUSY $<
100$~TeV, de désintégration du proton) sont soit différées soit non-falsifiables à court terme.

La Proposition~\ref{prop:c3-residue} conjointement avec le Théorème~\ref{thm:c3-uniqueness}
positionne Casimir comme test crucial expérimental le plus accessible de la classification PT
des nombres premiers.

\section{Conclusion}

Nous avons établi que le coefficient standard de la pression de Casimir entre plaques idéales,
$\pi^{2}/(240\,a^{4})$, admet une dérivation PT-native sans paramètre ajusté. La construction
comporte trois composantes~:
\begin{enumerate}
  \item Le préfacteur géométrique
  $\mathcal{G}_{3} = \Nspatial/(2\,(2\pi)^{\Ntrans})$, identifié entièrement par le décompte
  des dimensions spatiales en PT.
  \item Le produit eulérien $\zeta(4) = \prod_{p}(1-p^{-4})^{-1}$, décomposé selon la
  classification PT des nombres premiers en parité, actifs et d'écho.
  \item La troncature C3 à $\PTset$ prédit le coefficient standard à $131$~ppm près, le résidu
  étant rigoureusement la queue eulérienne d'écho $\prod_{p \geq 11}(1-p^{-4})^{-1}$.
\end{enumerate}

La construction est universelle pour les dimensions spatiales testées ($d \in \{1,\dots,5\}$),
la pureté de la prédiction PT croissant super-exponentiellement avec $d$. La position de
Casimir 3D dans le \emph{point d'équilibre expérimental} ($131$~ppm) en fait la sonde la plus
accessible de la classification PT.

Conjointement, l'identité fondamentale $\oint \mathrm{d}\IPT = 0$ exclut tout dispositif
d'extraction d'énergie nette à cavité statique~: la lecture PT renforce les contraintes
thermodynamiques en les réinterprétant comme contraintes arithmétiques exactes.

Une mesure différentielle de Casimir à mieux que $100$~ppm constituerait, suivant son résultat,
soit la première détection directe des nombres premiers d'écho comme structure physique
réelle, soit une falsification de la classification PT. Dans les deux cas, le résultat serait
décisif.

\section*{Remerciements}

L'auteur remercie la communauté open-source PT (PT\_CHEMISTRY, PT\_PHYSICS,
\url{https://github.com/Igrekess/PT_CHEMISTRY}) pour l'infrastructure computationnelle. Les
calculs numériques de cet article sont reproductibles \emph{via} les scripts
\path{run_p2_coefficient_routes.py} et \path{run_p2_dimensional_scan.py} du dépôt
\path{PT_EXPERIENCES/PT_CASIMIR/scripts/}.

\begin{thebibliography}{99}

\bibitem{Casimir1948}
H.~B.~G.~Casimir, ``On the attraction between two perfectly conducting plates,''
Proc.\ Kon.\ Nederland.\ Akad.\ Wetensch.\ \textbf{B51}, 793 (1948).

\bibitem{Lamoreaux1997}
S.~K.~Lamoreaux, ``Demonstration of the Casimir force in the 0.6 to 6~$\mu$m range,''
Phys.\ Rev.\ Lett.\ \textbf{78}, 5 (1997).

\bibitem{Mohideen1998}
U.~Mohideen et A.~Roy, ``Precision measurement of the Casimir force from 0.1 to 0.9~$\mu$m,''
Phys.\ Rev.\ Lett.\ \textbf{81}, 4549 (1998).

\bibitem{Bressi2002}
G.~Bressi, G.~Carugno, R.~Onofrio et G.~Ruoso,
``Measurement of the Casimir force between parallel metallic surfaces,''
Phys.\ Rev.\ Lett.\ \textbf{88}, 041804 (2002).

\bibitem{Decca2003}
R.~S.~Decca \emph{et al.},
``Measurement of the Casimir force between dissimilar metals,''
Phys.\ Rev.\ Lett.\ \textbf{91}, 050402 (2003).

\bibitem{Decca2007}
R.~S.~Decca \emph{et al.},
``Tests of new physics from precise measurements of the Casimir pressure between two
gold-coated plates,''
Phys.\ Rev.\ D \textbf{75}, 077101 (2007).

\bibitem{Lifshitz1956}
E.~M.~Lifshitz, ``The theory of molecular attractive forces between solids,''
Sov.\ Phys.\ JETP \textbf{2}, 73 (1956).

\bibitem{Mostepanenko2009}
M.~Bordag, G.~L.~Klimchitskaya, U.~Mohideen et V.~M.~Mostepanenko,
\emph{Advances in the Casimir Effect}, International Series of Monographs on Physics
\textbf{145}, Oxford University Press (2009).

\bibitem{Klimchitskaya2009}
G.~L.~Klimchitskaya, U.~Mohideen et V.~M.~Mostepanenko,
``The Casimir force between real materials: Experiment and theory,''
Rev.\ Mod.\ Phys.\ \textbf{81}, 1827 (2009).

\bibitem{Milton2001}
K.~A.~Milton, \emph{The Casimir Effect: Physical Manifestations of Zero-Point Energy},
World Scientific (2001).

\bibitem{Wilson2011}
C.~M.~Wilson, G.~Johansson, A.~Pourkabirian, M.~Simoen, J.~R.~Johansson, T.~Duty, F.~Nori et
P.~Delsing, ``Observation of the dynamical Casimir effect in a superconducting circuit,''
Nature \textbf{479}, 376 (2011).

\bibitem{Feynman1963}
R.~P.~Feynman, R.~B.~Leighton et M.~Sands,
\emph{The Feynman Lectures on Physics}, vol.~I, ch.~46 (Ratchet and pawl),
Addison-Wesley (1963).

\bibitem{Pinto1999}
F.~Pinto, ``Engine cycle of an optically controlled vacuum energy transducer,''
Phys.\ Rev.\ B \textbf{60}, 14740 (1999). [Le claim a été contesté~; voir aussi
brevet US 6\,477\,028.]

\bibitem{Senez2026}
Y.~Senez, \emph{Persistence Theory: A Complete Monograph}, monographie interne, mai 2026.
Disponible aux versions française et anglaise (\texttt{main\_fr.tex}, \texttt{main.tex}).

\bibitem{Senez2026RH}
Y.~Senez, \emph{Contributions de la Théorie de la Persistance à l'hypothèse de Riemann},
companion paper, 2026.

\bibitem{Senez2026HL}
Y.~Senez, \emph{La route PT vers Hardy-Littlewood}, companion paper, 2026.

\end{thebibliography}

\end{document}
